\newpage
\selectlanguage{greek}
\chapter*{Περίληψη}
\addcontentsline{toc}{chapter}{Περίληψη}

%Αυτή η διπλωματική εργασία παρουσιάζει το \textlatin{\textbf{S\textsuperscript{4}D} (Style and Structure Similarity for Site Domains)}, μια ερευνητική πλατφόρμα για την ανάλυση ομοιοτήτων μεταξύ ιστοσελίδων βάσει των οπτικών τους χαρακτηριστικών - \textlatin{Cascading Style Sheets (CSS)}. Το σύστημα επεξεργάζεται ιστοσελίδες αναλύοντας τη δομή του \textlatin{DOM} και τις ιδιότητες \textlatin{CSS}, συνδυάζοντάς τα σε μία ενιαία μορφή για συγκριτική αξιολόγηση.

%Οι χρήστες μπορούν να εισάγουν οποιονδήποτε αριθμό \textlatin{URLs} για ανάλυση. Κάθε σελίδα επεξεργάζεται ανεξάρτητα για την εξαγωγή δομικών και στυλιστικών χαρακτηριστικών. Κατά τη φάση υπολογισμού της ομοιότητας, οι συγκρίσεις πραγματοποιούνται σε επίπεδο \textlatin{domain}, ανεξάρτητα από τον τρόπο ομαδοποίησης των \textlatin{URLs}. Το σύστημα υποστηρίζει την ανάλυση σε διαφορετικά βάθη \textlatin{DOM} (επίπεδα 1 έως 10).

%Η ομαδοποίηση των στοιχείων πραγματοποιείται με τον αλγόριθμο \textlatin{HDBSCAN}, ενώ η ποιότητα των ομάδων αξιολογείται με το \textlatin{DBCV}. Τα απομονωμένα στοιχεία επανατοποθετούνται σε κατάλληλες ομάδες όταν είναι δυνατόν. Η πλατφόρμα επιτρέπει τη διεξαγωγή πολλαπλών πειραμάτων με διαφορετικά σύνολα \textlatin{URLs} και παραμέτρους, καλύπτοντας ποικίλες ερευνητικές ανάγκες.

%Το \textlatin{\textbf{S\textsuperscript{4}D}} υλοποιείται με \textlatin{React} για το περιβάλλον χρήστη, \textlatin{Flask} για τις διεπαφές \textlatin{API} και \textlatin{Python} για την επεξεργασία δεδομένων. Ενσωματώνει το \textlatin{Gnuplot} για παραγωγή γραφημάτων και αρχείων κατάλληλων για ακαδημαϊκή χρήση. Ένα αρχείο αποτελεσμάτων επιτρέπει τη σύγκριση διαφορετικών πειραμάτων και την παρακολούθηση της προόδου.

%Το σύστημα εξάγει \textlatin{HTML} σελίδες με χρωματική κωδικοποίηση που αναδεικνύουν τις δομικές ομοιότητες, καθώς και πίνακες συσχετίσεων βασισμένους σε υπολογισμούς συνημιτόνου. Οι χρονικές μετρήσεις συμβάλλουν στην αξιολόγηση της απόδοσης του συστήματος. Οι δοκιμές δείχνουν ότι το \textlatin{\textbf{S\textsuperscript{4}D}} συγκρίνει και ομαδοποιεί με ακρίβεια στοιχεία ιστοσελίδων σε διαφορετικά \textlatin{domains}, προσφέροντας ένα αξιόπιστο εργαλείο για δομική και οπτική ανάλυση.

%Η πλήρης υλοποίηση του συστήματος είναι διαθέσιμη στο: \textlatin{\url{https://github.com/pfavvatas/lib_url_to_img}}

Τα τελευταία χρόνια, οι ιστότοποι έχουν γίνει πιο πολύπλοκοι και ποικιλόμορφοι ως προς τον σχεδιασμό και τη δομή τους. Παρά την ποικιλομορφία αυτή, οι ιστότοποι που βρίσκονται στον ίδιο θεματικό τομέα ή εξυπηρετούν παρόμοιους σκοπούς συχνά μοιράζονται κοινά παρουσιαστικά και δομικά μοτίβα. Η αναγνώριση και η σύγκριση αυτών των μοτίβων σε διαφορετικούς τομείς ιστότοπων είναι πολύτιμη για την επίτευξη αποτελεσμάτων ως προς την ταξινόμηση ιστοσελίδων, την αυτοματοποιημένη ανάλυση διεπαφών χρήστη \textlatin{(UIs)} και την αποτύπωση των θεματικών τους τομέων. Ωστόσο, η δομή και η οπτική παρουσίαση των περιεχομένων ιστού είναι συχνά μη τυποποιημένη και ασυνεπής, αυξάνοντας τη δυσκολία τέτοιων αναλύσεων. Σε αυτήν την πτυχιακή εργασία, παρουσιάζεται το \textlatin{\textbf{S\textsuperscript{4}D}}, ένα σύστημα που αυτοματοποιεί τη διαδικασία εξαγωγής, μοντελοποίησης και σύγκρισης των παρουσιαστικών και δομικών ομοιοτήτων διαφορετικών ιστοτόπων. Το σύστημα αναλύει ένα σύνολο ιστοσελίδων μέσω τεχνικών αυτοματοποιημένης περιήγησης, μοντελοποιώντας τα \textlatin{CSS} χαρακτηριστικά των στοιχείων τους και οργανώνοντάς τα με τρόπο που επιτρέπει την περαιτέρω ανάλυση. Συγκεκριμένα, μετατρέποντας τα στοιχεία σελίδας σε διανυσματικές αναπαραστάσεις και εφαρμόζοντας τεχνικές μη-επιβλεπόμενης συσταδοποίησης, το σύστημα εντοπίζει ουσιαστικές συστάδες των στοιχείων βάσει των οπτικών τους χαρακτηριστικών. Χρησιμοποιώντας αυτήν την συσταδοποιημένη αναπαράσταση, το σύστημα μοντελοποιεί τα παραγόμενα δεδομένα για κάθε ιστότοπο, εξάγει μοτίβα διάταξης στοιχείων και υπολογίζει την ομοιότητα μεταξύ τους χρησιμοποιώντας διαφορετικές μετρικές απόστασης. Τέλος, το σύστημα παρουσιάζει έναν πίνακα ομοιότητας μεταξύ των ιστοτόπων, προσφέροντας πληροφορίες σχετικά με τη δομική και παρουσιαστική ομοιότητα των ιστοσελίδων τους.

%In recent years, websites have become increasingly complex and diverse in their design and structure. Despite this diversity, websites within the same domain or serving similar purposes often share common stylistic and structural patterns. Identifying and comparing these patterns across different site domains is valuable for tasks such as web classification, automated UI analysis, and domain fingerprinting. However, the structure and visual presentation of web content is often unstandardized and inconsistent, making such analysis a challenging task. This thesis introduces S4D, a system that automates the process of extracting, modeling, and comparing the style and structure of websites across different domains. The system parses a set of webpages using browser automation, models CSS-related properties of web elements, and structures them in a way that enables further analysis. By transforming page elements into vector representations and applying unsupervised clustering techniques, the system identifies meaningful groupings of visual and structural components. Using this clustered representation, the system performs data modeling for each domain, extracts high-level patterns of element arrangement, and computes inter-domain similarity using multiple distance metrics. Finally, the system presents a similarity matrix between domains, offering insights into the structural and stylistic resemblance of their webpages.